\begin{proposition}
\label{prop:ch28-curve-dual-energy}
\index{Dirichlet energy!on a metric graph}\index{Green function!energy of a measure}
Let \(X\) be a smooth projective curve under \eqref{eq:ch28-standing},
\(\Gamma\) the skeleton of a regular semistable model determining the
semipositive real form \(\omega\), and \(\mu_0=\MA_\omega(0)\), of
mass \(V=(\{\omega\})>0\).  Let \(\mu\) be any positive Radon measure
of mass \(V\) carried by \(\Gamma\), and let \(u\) be the continuous
BFJ solution constructed by \cref{lem:ch26-weighted-graph-real-bfj}.
With the residue-weighted energy \eqref{eq:ch28-dirichlet-form},
\begin{equation}
\label{eq:ch28-curve-dual-energy}
 E^\vee_\omega(\mu)
 =\tfrac12\sum_e c_e\int_e|u'|^2\,\rd s
 =\tfrac12D_\Gamma(u)
 =\tfrac12\int_\Gamma u\,\rd(\mu_0-\mu)
 \in[0,+\infty).
\end{equation}
In particular \(E^\vee_\omega(\mu)=0\) if and only if \(\mu=\mu_0\).
No atomicity or rationality of \(\mu\) is assumed.
\end{proposition}

\begin{proof}
The construction supplies an actual continuous BFJ-psh function with
\(\MA_\omega(u)=\mu\), for the full real-form scope of this chapter.
It is bounded, so it lies in \(\cE^1(X,\omega)\) by
\cref{thm:ch26-e1-structure}\,\cref{itm:ch26-e1:bounded}.
Therefore \cref{prop:ch28-dual-energy-of-ma} gives
\[
 E^\vee_\omega(\mu)=E_\omega(u)-\int u\,\rd\mu.
\]
The continuous case of \cref{prop:ch26-graph-energy} gives
\[
 E_\omega(u)=\int_\Gamma u\,\rd\mu_0-\tfrac12D_\Gamma(u).
\]
That continuous case is justified by positive real model approximants
whose derivatives converge strongly in the weighted edge \(L^2\) spaces;
uniform convergence alone does not imply convergence of their energies.
Finally \eqref{eq:ch26-weighted-green-ibp} reads
\(D_\Gamma(u)=\int_\Gamma u\,\rd(\mu_0-\mu)\).
Substitution gives \eqref{eq:ch28-curve-dual-energy}.
Every term is finite by the same lemma.
If the energy is zero, the derivative vanishes almost everywhere on every
edge, whose weight is positive.  Absolute continuity and connectedness
make \(u\) constant, so \(\mu=\mu_0+\Delta_\Gamma u=\mu_0\).
Conversely \(u=0\) solves the reference equation and has energy zero.
The statement is unchanged by adding a constant to \(u\), since the
two measures have equal mass.
\end{proof}

\begin{remark}[The classical energy of a measure]
\label{rem:ch28-classical-energy}
Formula \eqref{eq:ch28-curve-dual-energy} identifies the Legendre dual
energy with half the \emph{weighted} Dirichlet energy of the graph
potential relative to \(\mu_0\).  Choose \(q\in\Gamma\), use the
kernel \(g_q\) of \cref{lem:ch31-green-kernel}, and put
\(\nu=\mu-\mu_0\).  Up to a constant,
\(u(x)=\int_\Gamma g_q(x,y)\,\rd\nu(y)\), so
\[
 E^\vee_\omega(\mu)
 =-\tfrac12\iint_{\Gamma\times\Gamma}
          g_q(x,y)\,\rd\nu(x)\,\rd\nu(y).
\]
The double integral exists because the kernel is bounded and continuous
and the signed measure has finite total variation.  Its sign follows
from the outgoing-slope convention and the nonnegative weighted energy.
The zero mass of \(\nu\) removes dependence on the normalizing constant.
This is the curve quadratic form; in higher dimension the replacement is
the mixed-intersection theory of \cref{thm:ch26-mixed-energy-identities}.

For a nonsplit numerical check, take
\cref{ex:ch26-nonsplit-weighted-energy}, with
\(\mu_0=2\delta_{v_D}\), \(\mu=2\delta_{v_E}\), and endpoint
potential values \((0,-1)\).  Then
\[
 E^\vee_\omega(2\delta_{v_E})
 =E_\omega(u)-\int u\,\rd\mu=-1-(-2)=1
 =\tfrac12\cdot2.
\]
An unweighted derivative integral would give \(1/2\), contradicting
the variational value of this actual nef model metric.
\end{remark}
